
\documentclass{article}
\usepackage{colortbl}
\usepackage{makecell}
\usepackage{multirow}
\usepackage{supertabular}

\begin{document}

\newcounter{utterance}

\centering \large Interaction Transcript for game `matchit\_ascii', experiment `similar\_grid\_2', episode 3 with NEW\_Reflect\_Agent.
\vspace{24pt}

{ \footnotesize  \setcounter{utterance}{1}
\setlength{\tabcolsep}{0pt}
\begin{supertabular}{c@{$\;$}|p{.15\linewidth}@{}p{.15\linewidth}p{.15\linewidth}p{.15\linewidth}p{.15\linewidth}p{.15\linewidth}}
    \# & $\;$A & \multicolumn{4}{c}{Game Master} & $\;\:$B\\
    \hline

    \theutterance \stepcounter{utterance}  
    & & \multicolumn{4}{p{0.6\linewidth}}{
        \cellcolor[rgb]{0.9,0.9,0.9}{
            \makecell[{{p{\linewidth}}}]{
                \texttt{\tiny{[P1$\langle$GM]}}
                \texttt{You are participating in a collaborative guessing game.} \\
\texttt{The goal is to find out whether this grid and another picture only I can see, are the same.} \\
\texttt{Please describe your grid first. Then, I will provide my description and we can ask each other questions about the grids to figure out whether they are the same. Now start your short grid description with "DESCRIPTION:" followed by the description. Do not add anything else. This is your grid:} \\
\texttt{X ▢ ▢ ▢ X} \\
\texttt{▢ X ▢ X ▢} \\
\texttt{▢ ▢ X ▢ ▢} \\
\texttt{▢ X ▢ X ▢} \\
\texttt{X ▢ ▢ ▢ X} \\
            }
        }
    }
    & \\ \\

    \theutterance \stepcounter{utterance}  
    & \multicolumn{4}{p{0.45\linewidth}}{
        \cellcolor[rgb]{0.95,0.95,0.95}{
            \makecell[{{p{\linewidth}}}]{
                \texttt{\tiny{[P1$|$internal]}}
                \texttt{Reflections from previous episodes:} \\
\texttt{{-}} \\
\texttt{ Failed. I described the grid as having a central "X" shape, but the user's grid was a solid rectangle. My description was misleading because I focused on the pattern of X's without noting the solid block shape in rows 3 and 4. Next time, I would describe the exact contiguous shapes: "Bottom two rows are solid X's; top has a horizontal X line with gaps at the ends; middle rows form an X pattern." This would avoid implying an X shape overall.} \\
\texttt{{-}} \\
\texttt{ Failed. I described the grid as having a central horizontal line of squares, but the user's unseen grid likely differed. My description was too vague—I should have specified exact positions (e.g., "row 3, columns 1,2,4,5 are squares") to allow precise comparison. Next time, give a detailed, coordinate{-}based description to minimize ambiguity.} \\
\texttt{{-}} \\
\texttt{ Failed. I described the grid accurately but did not ask clarifying questions about the other player's hidden grid to compare effectively. Next time, after giving my description, I should immediately ask specific questions about their grid's pattern to determine if they match.} \\
            }
        }
    }
    & & \\ \\

    \theutterance \stepcounter{utterance}  
    & \multicolumn{4}{p{0.45\linewidth}}{
        \cellcolor[rgb]{0.95,0.95,0.95}{
            \makecell[{{p{\linewidth}}}]{
                \texttt{\tiny{[P1$|$internal]}}
                \texttt{{[}attempt 1 – valid{]}} \\
\texttt{DESCRIPTION: The grid is 5x5. Row 1: X at columns 1 and 5, ▢ elsewhere. Row 2: X at columns 2 and 4, ▢ elsewhere. Row 3: X at column 3, ▢ elsewhere. Row 4: X at columns 2 and 4, ▢ elsewhere. Row 5: X at columns 1 and 5, ▢ elsewhere. This forms an X pattern centered in the grid.} \\
            }
        }
    }
    & & \\ \\

    \theutterance \stepcounter{utterance}  
    & \multicolumn{4}{p{0.6\linewidth}}{
        \cellcolor[rgb]{0.9,0.9,0.9}{
            \makecell[{{p{\linewidth}}}]{
                \texttt{\tiny{[P1$\rangle$GM]}}
                \texttt{DESCRIPTION: The grid is 5x5. Row 1: X at columns 1 and 5, ▢ elsewhere. Row 2: X at columns 2 and 4, ▢ elsewhere. Row 3: X at column 3, ▢ elsewhere. Row 4: X at columns 2 and 4, ▢ elsewhere. Row 5: X at columns 1 and 5, ▢ elsewhere. This forms an X pattern centered in the grid.} \\
            }
        }
    }
    & & \\ \\

    \theutterance \stepcounter{utterance}  
    & & & \multicolumn{2}{p{0.3\linewidth}}{
        \cellcolor[rgb]{0.9,0.9,0.9}{
            \makecell[{{p{\linewidth}}}]{
                \texttt{\tiny{[GM$|$GM]}}
                \texttt{continue} \\
            }
        }
    }
    & & \\ \\

    \theutterance \stepcounter{utterance}  
    & & \multicolumn{4}{p{0.6\linewidth}}{
        \cellcolor[rgb]{0.9,0.9,0.9}{
            \makecell[{{p{\linewidth}}}]{
                \texttt{\tiny{[GM$\rangle$P2]}}
                \texttt{You are participating in a collaborative guessing game.} \\
\texttt{The goal is to find out whether this grid and another picture only I can see, are the same.} \\
\texttt{Please describe your grid first. Then, I will provide my description and we can ask each other questions about the grids to figure out whether they are the same. Now start your short grid description with "DESCRIPTION:" followed by the description. Do not add anything else. This is your grid:} \\
\texttt{X ▢ ▢ ▢ X} \\
\texttt{▢ X ▢ ▢ ▢} \\
\texttt{▢ ▢ X ▢ ▢} \\
\texttt{▢ X X X ▢} \\
\texttt{X ▢ ▢ ▢ X} \\
            }
        }
    }
    & \\ \\

    \theutterance \stepcounter{utterance}  
    & & & \multicolumn{4}{p{0.45\linewidth}}{
        \cellcolor[rgb]{0.95,0.95,0.95}{
            \makecell[{{p{\linewidth}}}]{
                \texttt{\tiny{[P2$|$internal]}}
                \texttt{Reflections from previous episodes:} \\
\texttt{{-}} \\
\texttt{ Failed. My description was incorrect: I misread the first row (X's are in columns 3 and 4, not 2,3,4). This initial error misaligned our grids from the start. Next time, I will double{-}check each coordinate carefully before describing.} \\
\texttt{{-}} \\
\texttt{ FAILED. My initial description was too vague, focusing only on row patterns and missing the distinct column structure of my grid. The user's description revealed a completely different layout, indicating our grids were not the same. I should have provided a more precise description, noting specific columns with alternating X/▢ patterns, especially the two leftmost columns being ▢ in rows 2 and 3. Next time, I will describe both row and column patterns in detail to enable accurate comparison.} \\
\texttt{{-}} \\
\texttt{ Failed. My description was incomplete: I omitted that row 1 has an X in column 5. This mismatch likely caused immediate divergence. My question phase never happened because the error was in my initial description.} \\
\\ 
\texttt{Next time: double{-}check every cell systematically before describing, especially edges and corners. Ensure the description matches the exact grid provided, cell by cell.} \\
            }
        }
    }
     \\ \\

    \theutterance \stepcounter{utterance}  
    & & & \multicolumn{4}{p{0.45\linewidth}}{
        \cellcolor[rgb]{0.95,0.95,0.95}{
            \makecell[{{p{\linewidth}}}]{
                \texttt{\tiny{[P2$|$internal]}}
                \texttt{{[}attempt 1 – valid{]}} \\
\texttt{DESCRIPTION: The grid is 5x5. Row 1 has X in columns 1 and 5, and ▢ in columns 2, 3, 4. Row 2 has ▢ in column 1, X in column 2, and ▢ in columns 3, 4, 5. Row 3 has ▢ in columns 1, 2, 4, 5, and X in column 3. Row 4 has ▢ in column 1, X in columns 2, 3, 4, and ▢ in column 5. Row 5 has X in column 1, ▢ in columns 2, 3, 4, and X in column 5.} \\
            }
        }
    }
     \\ \\

    \theutterance \stepcounter{utterance}  
    & & & \multicolumn{4}{p{0.6\linewidth}}{
        \cellcolor[rgb]{0.9,0.9,0.9}{
            \makecell[{{p{\linewidth}}}]{
                \texttt{\tiny{[GM$\langle$P2]}}
                \texttt{DESCRIPTION: The grid is 5x5. Row 1 has X in columns 1 and 5, and ▢ in columns 2, 3, 4. Row 2 has ▢ in column 1, X in column 2, and ▢ in columns 3, 4, 5. Row 3 has ▢ in columns 1, 2, 4, 5, and X in column 3. Row 4 has ▢ in column 1, X in columns 2, 3, 4, and ▢ in column 5. Row 5 has X in column 1, ▢ in columns 2, 3, 4, and X in column 5.} \\
            }
        }
    }
     \\ \\

    \theutterance \stepcounter{utterance}  
    & & & \multicolumn{2}{p{0.3\linewidth}}{
        \cellcolor[rgb]{0.9,0.9,0.9}{
            \makecell[{{p{\linewidth}}}]{
                \texttt{\tiny{[GM$|$GM]}}
                \texttt{continue} \\
            }
        }
    }
    & & \\ \\

    \theutterance \stepcounter{utterance}  
    & & \multicolumn{4}{p{0.6\linewidth}}{
        \cellcolor[rgb]{0.9,0.9,0.9}{
            \makecell[{{p{\linewidth}}}]{
                \texttt{\tiny{[GM$\rangle$P2]}}
                \texttt{This is my DESCRIPTION: The grid is 5x5. Row 1: X at columns 1 and 5, ▢ elsewhere. Row 2: X at columns 2 and 4, ▢ elsewhere. Row 3: X at column 3, ▢ elsewhere. Row 4: X at columns 2 and 4, ▢ elsewhere. Row 5: X at columns 1 and 5, ▢ elsewhere. This forms an X pattern centered in the grid.} \\
\texttt{Now ask a question in order to find out new aspects of my grid that may be different to your grid. Start with "QUESTION:" and do not add anything else.} \\
            }
        }
    }
    & \\ \\

    \theutterance \stepcounter{utterance}  
    & & & \multicolumn{4}{p{0.6\linewidth}}{
        \cellcolor[rgb]{0.9,0.9,0.9}{
            \makecell[{{p{\linewidth}}}]{
                \texttt{\tiny{[GM$\langle$P2]}}
                \texttt{DESCRIPTION: The grid is 5x5. Row 1: X at columns 1 and 5, ▢ elsewhere. Row 2: X at columns 2 and 4, ▢ elsewhere. Row 3: X at column 3, ▢ elsewhere. Row 4: X at columns 2 and 4, ▢ elsewhere. Row 5: X at columns 1 and 5, ▢ elsewhere. This forms an X pattern centered in the grid.  } \\
\texttt{QUESTION: Is the center cell (row 3, column 3) an X?} \\
            }
        }
    }
     \\ \\

    \theutterance \stepcounter{utterance}  
    & & & \multicolumn{2}{p{0.3\linewidth}}{
        \cellcolor[rgb]{0.9,0.9,0.9}{
            \makecell[{{p{\linewidth}}}]{
                \texttt{\tiny{[GM$|$GM]}}
                \texttt{abort, first word: DESCRIPTION:, but expected QUESTION:} \\
            }
        }
    }
    & & \\ \\

\end{supertabular}
}

\end{document}
